\documentclass[11pt]{article}
\usepackage[T1]{fontenc}
\usepackage{lmodern}
\usepackage[margin=1.05in]{geometry}
\usepackage{amsmath,amssymb,amsthm,mathrsfs}
\usepackage{microtype,booktabs}
\usepackage[colorlinks=true,linkcolor=blue,citecolor=blue,urlcolor=blue]{hyperref}
\hypersetup{pdftitle={Logarithmic lower bounds for the block complexity of algebraic numbers},
 pdfsubject={Diophantine approximation and expansions in integer bases}}
\numberwithin{equation}{section}
\newtheorem{theorem}{Theorem}[section]
\newtheorem{lemma}[theorem]{Lemma}
\newtheorem{proposition}[theorem]{Proposition}
\newtheorem{corollary}[theorem]{Corollary}
\theoremstyle{remark}
\newtheorem{remark}[theorem]{Remark}
\newcommand{\Q}{\mathbb Q}
\newcommand{\Z}{\mathbb Z}
\newcommand{\R}{\mathbb R}
\newcommand{\Qbar}{\overline{\mathbb Q}}
\newcommand{\eps}{\varepsilon}
\newcommand{\HH}{\mathcal H}
\DeclareMathOperator{\Gal}{Gal}
\title{Logarithmic lower bounds for the block complexity\\of algebraic numbers}
\author{}
\date{}
\begin{document}
\maketitle

\begin{abstract}
Let $b\ge2$ be an integer, let $\xi$ be a real algebraic irrational number,
and let $p_{\xi,b}(L)$ count the distinct blocks of length $L$ in the
fractional base-$b$ expansion of $\xi$. We prove that
\[
 \limsup_{L\to\infty}
 \frac{p_{\xi,b}(L)(\log\log L)^\gamma}{L(\log L)^{1/3}}
 =+\infty
 \qquad\text{for every }\gamma>1.
\]
In particular, $p_{\xi,b}(L)/(L(\log L)^\eta)$ is unbounded for every
real $\eta<1/3$.
The proof combines rational approximations obtained from repeated blocks
with the twisted-height interval theorem of Evertse and Ferretti.
A selection of the approximations has uniformly bounded intersection
with every proper linear subspace. This bound follows from a divisibility
condition and an algebraic norm estimate. For composite bases, the
approximating vectors need not be primitive; a direct use of the external
parameter in the twisted height avoids normalisation by their common
divisors.
\end{abstract}

\medskip
\noindent\textbf{Draft status and AI disclosure.}
The overall proof idea was generated by GPT, which was also used in drafting
and internal checks. This version is shared for discussion and has not
received independent mathematical peer review.

\section{Introduction}\label{sec:intro}

Let $b\ge2$ be an integer and write the fractional part of a real irrational
number $\xi$ as
\[
 \{\xi\}=\sum_{i\ge1}d_i b^{-i},\qquad d_i\in\{0,\ldots,b-1\}.
\]
Its block complexity is
\[
 p_{\xi,b}(L)
 =\#\{d_{j+1}\cdots d_{j+L}:j\ge0\},\qquad L\ge1.
\]
Normality to base $b$ implies $p_{\xi,b}(L)=b^L$ for every $L$.
Much weaker lower bounds are already of interest when $\xi$ is algebraic.

Adamczewski and Bugeaud \cite[Theorem~1]{AB07} proved that, for every
algebraic irrational $\xi$ and every integer $b\ge2$,
\[
 \lim_{L\to\infty}\frac{p_{\xi,b}(L)}{L}=+\infty.
\]
Bugeaud and Evertse \cite[Theorem~2.1]{BE08} established the logarithmic
refinement
\[
 \limsup_{L\to\infty}
 \frac{p_{\xi,b}(L)}{L(\log L)^\eta}=+\infty
 \qquad(\eta<1/11).
\]
The result proved here is the following.

\begin{theorem}\label{thm:main}
For every fixed integer $b\ge2$, every real algebraic irrational number
$\xi$, and every real $\eta<1/3$, one has
\begin{equation}\label{eq:main}
 \limsup_{L\to\infty}
 \frac{p_{\xi,b}(L)}{L(\log L)^\eta}=+\infty.
\end{equation}
Equivalently, for every $C>0$ there are arbitrarily large integers $L$
such that $p_{\xi,b}(L)>CL(\log L)^\eta$.
\end{theorem}

Retaining a secondary logarithmic factor gives the following refinement.

\begin{theorem}\label{thm:endpoint}
For every fixed integer $b\ge2$, every real algebraic irrational number
$\xi$, and every real $\gamma>1$, one has
\begin{equation}\label{eq:endpoint}
 \limsup_{L\to\infty}
 \frac{p_{\xi,b}(L)(\log\log L)^\gamma}
      {L(\log L)^{1/3}}=+\infty.
\end{equation}
\end{theorem}

The logarithms in \eqref{eq:endpoint} are considered only for sufficiently
large $L$. This result implies Theorem~\ref{thm:main}, since
$(\log L)^{1/3-\eta}/(\log\log L)^\gamma\to\infty$ when $\eta<1/3$.
Both conclusions are limsup statements. The unweighted endpoint
$\limsup p_{\xi,b}(L)/(L(\log L)^{1/3})=+\infty$ and analogous lower
bounds at every sufficiently large length remain open here, as do
positive entropy, disjunctivity, and normality.
Bugeaud and Evertse also note a secondary logarithmic refinement at
their exponent $1/11$ in \cite[equation (2.4)]{BE08}.

We follow the repetition construction underlying \cite[\S6]{BE08}.
An eventual bound $p_{\xi,b}(L)\ll L(\log L)^u$ supplies approximations
\[
 |(b^t-b^r)\xi-a|\le b^{-v},\qquad
 v\gg \ell(\log\ell)^{-u},\qquad 0\le r<t\le\ell.
\]
Choosing a shortest preperiod gives $b\nmid a$ whenever $r>0$.
We then select a sequence for which the period lengths $s=t-r$
are distinct and the exponents $t$ more than double.
The selected vectors $(b^t,b^r,a)$ have uniformly bounded intersection
with every proper subspace. The bound depends only on $b$ and $\xi$,
and not on the approximation exponent; see Proposition~\ref{prop:occupancy}.

This intersection bound replaces the further application of a
two-dimensional approximation theorem used in \cite[\S7]{BE08}.
We combine it with the interval form of the quantitative Subspace
Theorem in \cite[Theorem~2.3]{EF13}. The uniform intersection bound
controls the exceptional subspace independently of the changing weights.
The remaining count has the following three factors, shown by their
orders of magnitude:
\[
 \underbrace{\eps^{-1}}_{\text{ratio intervals}}
 \underbrace{\eps^{-2}\log(2/\eps)}_{\text{EF intervals}}
 \underbrace{\log(2/\eps)}_{\text{points per interval}}
 =\eps^{-3}(\log(2/\eps))^2.
\]
The exceptional subspaces contribute only $O_{b,\xi}(\eps^{-1})$
additional points. Thus substituting a sharper quantitative theorem
alone is not the whole argument: the uniform subspace bound and the
selection of period lengths are also used.
Together with the growth of the selected prefix lengths, this gives
the range $u<1/3$ and the refinement in Theorem~\ref{thm:endpoint}.
Proposition~\ref{prop:bins} explains the cost of the ratio partition
within the fixed-weight estimates used here.

All arguments below apply to each fixed integer base. In composite bases,
$b\nmid a$ does not imply $\gcd(a,b)=1$. Consequently, the size $b^t$
cannot simply be identified with the projective height of the vector.
Section~\ref{sec:conversion} gives the local inequalities and their
conversion to twisted height with an external parameter.
Appendix~\ref{app:covering} records the weaker exponents obtained from
covering versions of the Subspace Theorem.

Throughout, logarithms are natural unless a base is specified.
Subtracting the integer part, we assume from now on that $0<\xi<1$,
and put $d=[\Q(\xi):\Q]$. For $x\in\Z^3\setminus\{0\}$, its absolute
projective height is
\[
 H(x)=\frac{\max_i|x_i|}{\gcd(x_1,x_2,x_3)}.
\]
Implied constants may depend on parameters explicitly fixed in the
argument. They are independent of the varying approximation exponent,
the chosen interval of ratios, the selected sequence length, and the
exceptional subspace.

\section{Repeated blocks and rational approximations}\label{sec:repetitions}

Fix $u>0$ and $\gamma\ge0$, and define, for real $x>e$,
\begin{equation}\label{eq:envelope}
 F(x)=\frac{(\log x)^u}{(\log\log x)^\gamma}.
\end{equation}
The case $\gamma=0$ gives the original logarithmic bound.

\begin{lemma}\label{lem:envelope}
The function $F$ is eventually increasing, tends to infinity, and
satisfies $F(x)=o(x)$. For every fixed $D>0$,
\begin{equation}\label{eq:slow}
 \frac{F(\lceil DxF(x)\rceil)}{F(x)}\longrightarrow1
 \qquad(x\to\infty).
\end{equation}
\end{lemma}
\begin{proof}
For sufficiently large $x$,
\[
 \frac{d}{dx}\log F(x)
 =\frac{u-\gamma/\log\log x}{x\log x}>0.
\]
Also $\log F(x)=u\log\log x-\gamma\log\log\log x$ tends to infinity
and is $o(\log x)$, which proves the asserted growth properties.
Writing $x^*=\lceil DxF(x)\rceil$, we have
$\log x^*=\log x+O_{D,u,\gamma}(\log\log x)$.
Thus both $\log x^*/\log x$ and
$\log\log x^*/\log\log x$ tend to $1$. Substituting these two ratios
in \eqref{eq:envelope} proves \eqref{eq:slow}.
\end{proof}

Suppose that, for some fixed $C\ge1$ and sufficiently large $k_0$,
\begin{equation}\label{eq:LC}
 p_{\xi,b}(k)\le CkF(k)\qquad(k\ge k_0).
\end{equation}
Increase $k_0$ so that $F$ is increasing, $F\ge1$, and
$\log\log x\ge1$ for $x\ge k_0$.
We use juxtaposition for concatenation of finite words; empty words are
allowed except where a positive length is specified.
Let
\[
 A_j=\sum_{i=1}^j d_i b^{j-i}\quad(j\ge1),\qquad A_0=0.
\]

\begin{lemma}\label{lem:repeat}
For every sufficiently large integer $\ell$, there is a decomposition
\begin{equation}\label{eq:repeat}
 d_1\cdots d_\ell=UVWVX,\qquad
 v:=|V|\ge a_0\ell/F(\ell),\qquad a_0=(24C)^{-1}.
\end{equation}
\end{lemma}
\begin{proof}
Put $k=\lfloor\ell/(4CF(\ell))\rfloor$.
For large $\ell$, one has $k\ge k_0$ and
\[
 p_{\xi,b}(k)\le CkF(\ell)\le\ell/4,
\]
using the monotonicity of $F$.
There are $\ell-k+1>\ell/4$ starting positions for length-$k$ blocks in
the prefix. Two of these blocks coincide. Let $D$ be the distance
between their starting positions.

If $D\ge k$, the two occurrences are disjoint.
If $D<k$, the substring of length $k+D$ containing them has period $D$.
Its first two consecutive blocks of length
\[
 v'=D\left\lfloor\frac{k+D}{2D}\right\rfloor
\]
are equal, since $D$ divides $v'$ and $2v'\le k+D$.
Moreover $v'\ge k/3$. Indeed, if $D\ge k/3$ then $v'\ge D$;
otherwise $v'>(k-D)/2>k/3$.
In either case there is a disjoint repeated block of length at least
$k/3$. For large $\ell$, the bound
$k\ge\ell/(8CF(\ell))$ proves \eqref{eq:repeat}.
\end{proof}

For each such $\ell$, choose a decomposition \eqref{eq:repeat} with
$|V|$ maximal, and among these choose one with $|U|$ minimal.
Any remaining choice is fixed arbitrarily. Write
\[
 r=|U|,\qquad s=|VW|,\qquad t=r+s,\qquad
 a=A_t-A_r,\qquad q=b^t-b^r.
\]

\begin{lemma}\label{lem:approx}
These integers satisfy
\begin{equation}\label{eq:structure}
 s\ge v,\qquad 0\le r<t\le\ell,\qquad
 0\le a\le q<b^t,\qquad |q\xi-a|\le b^{-v}.
\end{equation}
If $r>0$, then $b\nmid a$.
\end{lemma}
\begin{proof}
Suppose first that $r>0$ and $d_r=d_t$.
Moving both occurrences of $V$ one position to the left produces two
equal blocks of the same length and with the same separation.
They remain disjoint, and their first starting position is earlier.
This contradicts the minimality of $r$. Hence $d_r\ne d_t$, and
\[
 a=A_t-A_r\equiv d_t-d_r\not\equiv0\pmod b.
\]

If $P$ is the base-$b$ value of the word $VW$, then
$A_t=b^sA_r+P$, and therefore
\begin{equation}\label{eq:periodic}
 \frac aq
 =\frac{A_r}{b^r}+\frac{P}{b^r(b^s-1)}.
\end{equation}
The right side is the value of the infinite word $U(VW)^\infty$.
This word and the expansion of $\xi$ have the same first $t+v$
digits. Values with a given prefix of length $t+v$ lie in a closed
interval of length $b^{-(t+v)}$. This observation also covers a
periodic tail consisting entirely of the digit $b-1$. Thus
\[
 |\xi-a/q|\le b^{-(t+v)},\qquad |q\xi-a|\le b^{-v}.
\]
Finally $0\le A_r\le b^r-1$ and $0\le P\le b^s-1$ give
\[
 0\le a=(b^s-1)A_r+P\le b^r(b^s-1)=q<b^t.
\]
The remaining inequalities follow from the word lengths.
\end{proof}

\section{Selection of the approximations}\label{sec:selection}

Attach subscripts $\ell$ to the quantities constructed in
Section~\ref{sec:repetitions}.

\begin{lemma}\label{lem:selection}
There is an increasing sequence of admissible prefix lengths $\ell_j$
such that, on writing $r_j=r_{\ell_j}$ and similarly for $s_j,t_j,v_j,a_j$,
\begin{equation}\label{eq:selection}
 s_{j+1}>2t_j,\qquad t_{j+1}>2t_j,\qquad s_{j+1}>s_j.
\end{equation}
For constants $D_1,A>0$ independent of $j$, one also has
\begin{equation}\label{eq:growth}
 \ell_{j+1}\le D_1\ell_jF(\ell_j),\qquad
 \log\ell_j\le Aj\log(2j).
\end{equation}
\end{lemma}
\begin{proof}
Choose $\ell_1$ sufficiently large. Given $\ell_j$, let $\ell_{j+1}$
be the least integer $\ell>\ell_j$ for which $s_\ell>2t_{\ell_j}$.
Such an integer exists because
$s_\ell\ge a_0\ell/F(\ell)\to\infty$.
Since $t_j\ge s_j$, this rule implies all of \eqref{eq:selection}.

Fix $D>4/a_0$ and consider
$\ell^*=\lceil D\ell_jF(\ell_j)\rceil$.
By Lemma~\ref{lem:envelope}, $F(\ell^*)/F(\ell_j)\to1$.
Consequently, for all sufficiently large $\ell_j$,
\[
 s_{\ell^*}
 \ge a_0\ell^*/F(\ell^*)
 \ge \frac{a_0D}{2}\ell_j
 >2t_j.
\]
Thus $\ell_{j+1}\le\ell^*$, proving the first estimate in
\eqref{eq:growth} after increasing a fixed constant.

Since $F(x)\le(\log x)^u$ on the range in question, for
$y_j=\log\ell_j$ this estimate implies
$y_{j+1}\le y_j+u\log y_j+O(1)$.
Choose $A$ large enough to cover $y_1$ and to ensure, for every $j\ge1$,
\[
 u\log\bigl(Aj\log(2j)\bigr)+O(1)\le A\log(2j).
\]
This is possible since the left side grows logarithmically in both
$j$ and $A$, whereas the coefficient on the right is linear in $A$.
Induction, using
\[
 (j+1)\log(2j+2)-j\log(2j)\ge\log(2j),
\]
now gives $y_j\le Aj\log(2j)$.
The constants may depend on $C,u,\gamma,k_0$ and the initial choice $\ell_1$.
\end{proof}

Since $t_j\ge2^{j-1}t_1$ and $\ell_j\ge t_j$, the same sequence also
satisfies, for all sufficiently large $j$,
\begin{equation}\label{eq:loggrowth}
 j\ll\log\ell_j\ll j\log(2j),\qquad
 \log\log\ell_j\asymp\log(2j).
\end{equation}
The two-sided comparison will be used when a logarithmic factor has
a negative exponent.

Set
\begin{equation}\label{eq:points}
 x_j=(b^{t_j},b^{r_j},a_j),\qquad \Psi_j=b^{t_j}.
\end{equation}
In particular $\Psi_j=\max_i|(x_j)_i|$. If
$g_j=\gcd(b^{r_j},a_j)$, then
\begin{equation}\label{eq:ordinaryheight}
 H(x_j)=b^{t_j}/g_j,\qquad
 b^{s_j}\le H(x_j)\le b^{t_j}.
\end{equation}
Thus the stronger selection of the periods also implies
$H(x_{j+1})>H(x_j)^2$. The argument below, however, uses $\Psi_j$
as a size parameter without identifying it with $H(x_j)$.

\section{Uniform intersection with linear subspaces}\label{sec:occupancy}

The next proposition isolates the arithmetic property of the selected
sequence that will be used in the counting argument.

\begin{proposition}\label{prop:occupancy}
Fix $b\ge2$ and a real algebraic irrational $\xi\in(0,1)$.
Suppose that integer vectors
$x_j=(b^{t_j},b^{r_j},a_j)$ satisfy
\begin{align}
 &0\le r_j<t_j,\qquad 0\le a_j<b^{t_j},\qquad
 r_j>0\ \Longrightarrow\ b\nmid a_j,\label{eq:occ1}\\
 &|(b^{t_j}-b^{r_j})\xi-a_j|\le1,\label{eq:occ2}\\
 &t_{j+1}>2t_j,\qquad
 s_j:=t_j-r_j\text{ are pairwise distinct}.\label{eq:occ3}
\end{align}
There is a constant $B_{b,\xi}$ such that every proper linear subspace
of $\Qbar^3$ contains at most $B_{b,\xi}$ of the vectors $x_j$.
The constant is independent of the subspace and of the sequence.
\end{proposition}
\begin{proof}
Distinct vectors are not proportional. Indeed, if $t'>t$ and
$x'=\lambda x$, the first coordinates give $\lambda=b^{t'-t}$,
and the second give $r'=r+t'-t>0$. The third coordinate then gives
$b\mid a'$, contradicting \eqref{eq:occ1}.

Consider a rational plane with primitive integer normal vector
$z=(z_1,z_2,z_3)\ne0$, and put $Z=\max_i|z_i|$.
If the plane contains at least two selected vectors, let $x_i,x_j$,
$i<j$, be its first two. Their cross product is a nonzero integer
multiple of $z$. Since all coordinates of $x_h$ have absolute value
at most $b^{t_h}$, we have
\begin{equation}\label{eq:normal}
 Z\le2b^{t_i+t_j}\le2b^{2t_j}.
\end{equation}

If $z_3=0$, the equation of the plane is
$z_1b^t+z_2b^r=0$. When it has a selected solution, both $z_1,z_2$
are nonzero and $b^{t-r}=-z_2/z_1$ is fixed. The distinctness of the
$s_j$ therefore permits at most one selected vector.

Suppose that $z_3\ne0$, and write
$b=\prod_{p\mid b}p^{e_p}$. At a point with $r>0$, choose a prime
$p\mid b$ such that $v_p(a)<e_p$, as allowed by $b\nmid a$.
The plane equation implies $b^r\mid z_3a$, so
\[
 e_pr\le v_p(z_3)+v_p(a)\le v_p(z_3)+e_p-1.
\]
It follows, with $\kappa_b=\min_{p\mid b}e_p\log p>0$, that
\begin{equation}\label{eq:preperiodbound}
 r\le\frac{\log Z}{\kappa_b}+1.
\end{equation}
This also holds for $r=0$. The prime used in this argument is allowed
to depend on the point.

Let $A_d>0$ be the leading coefficient of the primitive minimal
polynomial of $\xi$, and let $\xi=\xi_1,\xi_2,\ldots,\xi_d$ be its
conjugates. Since $A_d\xi$ is an algebraic integer,
$A_d^dN_{\Q(\xi)/\Q}(z_1+z_3\xi)$ is a nonzero integer.
The other factors in this norm have absolute value at most
$Z(1+|\xi_h|)$. Hence
\begin{equation}\label{eq:norm}
 |z_1+z_3\xi|\ge c_\xi Z^{-(d-1)},\qquad
 c_\xi=\left(A_d^d\prod_{h=2}^d(1+|\xi_h|)\right)^{-1}>0.
\end{equation}
Write $a=\xi(b^t-b^r)+E$, where $|E|\le1$.
Substitution into the plane gives
\[
 (z_1+z_3\xi)b^t+(z_2-z_3\xi)b^r+z_3E=0.
\]
Using \eqref{eq:norm} and $b^r\ge1$, we obtain
\begin{equation}\label{eq:periodbound}
 b^{t-r}\le c_\xi^{-1}(2+\xi)Z^d.
\end{equation}
Together, \eqref{eq:preperiodbound} and \eqref{eq:periodbound} imply
\[
 t\le\left(\kappa_b^{-1}+\frac d{\log b}\right)\log Z
       +O_{b,\xi}(1).
\]
By \eqref{eq:normal}, every selected vector in this plane consequently
satisfies $t\le C_{b,\xi}t_j$ for a fixed $C_{b,\xi}\ge1$.
The doubling condition permits only a bounded number of indices after
$j$ with this property. There is just one earlier selected vector in
the plane, by the choice of $i,j$. This proves a uniform bound for
rational planes and therefore for proper rational subspaces.

Finally, if $T\subsetneq\Qbar^3$ is linear, then
$\dim_\Q(T\cap\Q^3)\le2$: otherwise three rationally independent
rational vectors in $T$ would have nonzero determinant.
If this dimension is $2$, the rational-plane bound applies.
If it is $1$, non-proportionality permits at most one selected vector;
if it is $0$, there is none, since each $x_j$ is nonzero.
Taking $B_{b,\xi}\ge1$ proves the assertion in every case.
\end{proof}

In particular, $B_{b,\xi}$ is independent of the plane, the chosen
sequence, its length $N$, the approximation parameter $\eps$, and
the ratio interval $I_l$. The exceptional subspace in the application
below may change with $N$ and $l$; this uniformity permits the same
bound to be used for all of them.

\begin{corollary}\label{cor:four}
For $b=2$ and $\xi=\sqrt2-1$, the constant in
Proposition~\ref{prop:occupancy} may be taken to be $4$.
\end{corollary}
\begin{proof}
Use the notation of the proof of Proposition~\ref{prop:occupancy}.
Only the case $z_3\ne0$ requires consideration.
If $r>0$, then $a$ is odd and the plane equation gives $2^r\mid z_3$.
Thus $2^r\le Z$, also for $r=0$.
The conjugate of $\xi$ is $-\sqrt2-1$, so \eqref{eq:norm} gives
$|z_1+z_3\xi|\ge((2+\sqrt2)Z)^{-1}$.
The same plane equation now yields
\[
 2^{t-r}
 \le(2+\sqrt2)(\sqrt2+1)Z^2
 =(4+3\sqrt2)Z^2<16Z^2.
\]
Consequently $t<3\log_2 Z+4\le6t_j+7$, where $x_i,x_j$ are the first
two selected vectors in the plane. Since the $t_h$ are integers,
$t_{h+1}\ge2t_h+1$, and hence
$t_{j+3}\ge8t_j+7>6t_j+7$.
There are no selected vectors of the plane with index at least $j+3$.
At most $x_i,x_j,x_{j+1},x_{j+2}$ can occur.
\end{proof}

\begin{remark}
The hypotheses of Proposition~\ref{prop:occupancy} require only the
bounded error \eqref{eq:occ2}. In particular, its constant remains
uniform as the approximation exponent tends to zero.
The fractional expansions of $\sqrt2$ and $\sqrt2-1$ coincide, so
Corollary~\ref{cor:four} applies to the approximations used for
$p_{\sqrt2,2}$.
\end{remark}

\section{A twisted-height interval theorem}\label{sec:EF}

We use the following immediate consequence of
\cite[Theorem~2.3]{EF13} in dimension three.
Let $K$ be a number field and $M_K$ its set of places.
Use absolute values $\|\cdot\|_v$ normalised for the product formula.
For every $v\in M_K$, let $L_{1v},L_{2v},L_{3v}$ be linearly independent
linear forms in $K[X_1,X_2,X_3]$. Assume that they are $X_1,X_2,X_3$
at all but finitely many places, and that at most $R$ distinct forms
occur altogether. Let real weights $c_{iv}$ have finite support in $v$
and satisfy
\begin{equation}\label{eq:weightconditions}
 \sum_{i=1}^3 c_{iv}=0\quad(v\in M_K),\qquad
 \sum_{v\in M_K}\max_{1\le i\le3}c_{iv}\le1.
\end{equation}
For $Q\ge1$ and $x\in K^3$, define
\[
 \HH_{\mathcal L,c,Q}(x)
 =\prod_{v\in M_K}\max_{1\le i\le3}
      \bigl(\|L_{iv}(x)\|_v Q^{-c_{iv}}\bigr),
 \qquad \mathcal L=(L_{iv}).
\]
The normalised extension of the forms and weights defines the same
height on $\Qbar^3$; only rational vectors will be needed here.
If $F_1,\ldots,F_R$ is the list of distinct forms, with repetitions
added if needed, put
\begin{align*}
 \Delta_{\mathcal L}
 &=\prod_{v\in M_K}\|\det(L_{1v},L_{2v},L_{3v})\|_v,\\
 H_{\mathcal L}
 &=\prod_{v\in M_K}
   \max_{1\le i_1,i_2,i_3\le R}
   \|\det(F_{i_1},F_{i_2},F_{i_3})\|_v.
\end{align*}

\begin{corollary}[of Evertse--Ferretti's interval theorem]\label{thm:intervals}
Under these hypotheses, let $0<\theta\le1$ and set
\[
 m=\left\lfloor10^5\,2^6\,3^{10}\theta^{-2}
                   \log(3\theta^{-1}R)\right\rfloor,\qquad
 \omega=\theta^{-1}\log(3R),\qquad
 C_0=\max\bigl(H_{\mathcal L}^{1/R},3^{1/\theta}\bigr).
\]
There is a proper linear subspace $T\subsetneq\Qbar^3$, defined over
$K$, and there are real numbers $Q_1,\ldots,Q_m\ge C_0$ such that,
for every $Q\ge1$ and $x\in\Qbar^3\setminus T$,
\[
 \HH_{\mathcal L,c,Q}(x)\le\Delta_{\mathcal L}^{1/3}Q^{-\theta}
 \quad\Longrightarrow\quad
 Q\in[1,C_0)\cup\bigcup_{h=1}^m[Q_h,Q_h^\omega).
\]
\end{corollary}
\begin{proof}
For fixed $Q$, the existence of an $x\notin T$ satisfying the displayed
inequality says exactly that its entire solution set is not contained
in $T$. Apply \cite[Theorem~2.3, equations (2.24)--(2.25)]{EF13}.
\end{proof}

The parameter $Q$ in this corollary is external; it need not equal
$H(x)$. The weights are fixed when the theorem is applied.
Both the exceptional subspace and the interval endpoints may depend
on these weights. The number of intervals is
$O_R(\theta^{-2}\log(2/\theta))$.

\section{Local inequalities and the height parameter}\label{sec:conversion}

Return to the sequence \eqref{eq:points} obtained under \eqref{eq:LC}.
We keep $\Psi_j=b^{t_j}$ as a size parameter; it is not identified
with $H(x_j)$. The parameter used in Corollary~\ref{thm:intervals}
will be $Q_j=\Psi_j^{1+\delta/3}$, so common divisors of the integer
coordinates require no additional partition.
Fix a sufficiently large even integer $N$, and put
\begin{equation}\label{eq:eps}
 \eps=a_0/F(\ell_N),\qquad
 k=\lceil2/\eps\rceil,\qquad
 \delta=\eps-1/k,\qquad
 \theta=\frac{\delta}{3+\delta}.
\end{equation}
For large $N$, one has
\begin{equation}\label{eq:exponentbounds}
 0<\eps\le1,\qquad \eps/2\le\delta\le\eps,\qquad
 \eps/8\le\theta\le\eps/3.
\end{equation}
For every $j\le N$, Lemmas~\ref{lem:repeat} and \ref{lem:approx} give
\begin{equation}\label{eq:epsilonapprox}
 v_j\ge a_0\ell_j/F(\ell_j)\ge\eps t_j,\qquad
 |(b^{t_j}-b^{r_j})\xi-a_j|\le\Psi_j^{-\eps}.
\end{equation}

Partition $[0,1)$ into the $k$ intervals
$I_l=[l/k,(l+1)/k)$, $0\le l<k$, and assign a point to $I_l$ when
$\rho_j=r_j/t_j\in I_l$.
Fix one such interval.
For $p\mid b$, put
\[
 w_p=\frac{e_p\log p}{\log b},\qquad
 \sum_{p\mid b}w_p=1,\qquad
 S=\{\infty\}\cup\{p:p\mid b\}.
\]
Here $|p|_p=p^{-1}$ and $|\cdot|_\infty$ is the usual real absolute
value. Choose an extension of each to $\Qbar$.
At the rational places in $S$, consider the forms and exponents
\begin{equation}\label{eq:raw}
\begin{array}{c|ccc|ccc}
 &L_{1,p}&L_{2,p}&L_{3,p}&e_{1,p}&e_{2,p}&e_{3,p}\\ \hline
 \infty&X_1&X_2&-\xi X_1+\xi X_2+X_3
       &1&(l+1)/k&-\eps\\
 p\mid b&X_1&X_2&X_3&-w_p&-w_pl/k&0 .
\end{array}
\end{equation}
Every selected vector assigned to $I_l$ satisfies
\begin{equation}\label{eq:rawineq}
 |L_{i,p}(x_j)|_p\le\Psi_j^{e_{i,p}}
 \qquad(p\in S,\ 1\le i\le3).
\end{equation}
At infinity this follows from the upper endpoint of $I_l$ and
\eqref{eq:epsilonapprox}. At $p\mid b$ it follows from
\[
 |b^{t_j}|_p=\Psi_j^{-w_p},\qquad
 |b^{r_j}|_p=\Psi_j^{-w_p\rho_j}\le\Psi_j^{-w_pl/k},
 \qquad |a_j|_p\le1.
\]
These are inequalities for integer coordinates without primitive
normalisation. Their exponents satisfy
\begin{equation}\label{eq:rawsums}
 \sum_{p\in S}\sum_i e_{i,p}=-\delta,\qquad
 \sum_{p\in S}\max_i e_{i,p}=1.
\end{equation}

For completeness, we give the conversion to twisted height, following
the construction in \cite[\S5, equations (5.14)--(5.17)]{BE08}.
Choose a finite Galois extension $K/\Q$ containing $\xi$ and its
conjugates. The previously fixed extensions $|\cdot|_p$ on $\Qbar$
serve as reference absolute values. For $v$ above a rational place $p$, choose
$\sigma_v\in\Gal(K/\Q)$ such that
\begin{equation}\label{eq:place}
 \|z\|_v=|\sigma_v(z)|_p^{d_v}\quad(z\in K),\qquad
 d_v=\frac{[K_v:\Q_p]}{[K:\Q]},\qquad
 \sum_{v\mid p}d_v=1.
\end{equation}
Such a choice exists because the Galois group acts transitively on
the places of $K$ above $p$. Make these choices once, independently
of the points and their parameters.
At infinity $\Q_\infty=\R$, with the corresponding real or complex
local degree. This is the normalisation of
\cite[\S1.1 and \S2.1, equations (2.1)--(2.3)]{EF13};
in particular a complex place has local degree $2$.

\begin{lemma}[Transport of a local form]\label{lem:transport}
For $v\mid p$, every linear form $L\in K[X_1,X_2,X_3]$ and every
$x\in\Q^3$ satisfy
\begin{equation}\label{eq:transport}
 \|\sigma_v^{-1}(L)(x)\|_v=|L(x)|_p^{d_v}.
\end{equation}
Here $\sigma_v^{-1}$ acts on the coefficients of the form.
\end{lemma}
\begin{proof}
The automorphism $\sigma_v$ fixes each coordinate of $x$. Thus
\[
 \sigma_v\bigl(\sigma_v^{-1}(L)(x)\bigr)
 =L(\sigma_v(x))=L(x).
\]
Applying \eqref{eq:place} proves \eqref{eq:transport}.
\end{proof}

At places above $S$ set
\[
 \widetilde L_{iv}=\sigma_v^{-1}(L_{i,p}),\qquad
 \widetilde e_{iv}=d_v e_{i,p};
\]
at all other places set $\widetilde L_{iv}=X_i$ and
$\widetilde e_{iv}=0$. For $v\mid p$ with $p\in S$,
Lemma~\ref{lem:transport} and \eqref{eq:rawineq} give the exact comparison
\[
 \|\widetilde L_{iv}(x_j)\|_v
 =|L_{i,p}(x_j)|_p^{d_v}
 \le\Psi_j^{d_ve_{i,p}}
 =\Psi_j^{\widetilde e_{iv}}.
\]
Together with integrality at places not above $S$, this proves
\begin{equation}\label{eq:lift}
 \|\widetilde L_{iv}(x_j)\|_v\le\Psi_j^{\widetilde e_{iv}}
 \qquad(v\in M_K).
\end{equation}

\begin{remark}\label{rem:twoembeddings}
For example, take $K=\Q(\sqrt2)$, $\xi=\sqrt2-1$, and
$\xi'=-\sqrt2-1$. Let $\tau(\sqrt2)=-\sqrt2$, and let $v_+,v_-$
be the two real places. With the positive embedding as reference,
\[
 \|z\|_{v_+}=|z|^{1/2},\qquad
 \|z\|_{v_-}=|\tau(z)|^{1/2}.
\]
For $q=b^t-b^r$ and $x=(b^t,b^r,a)$, the forms chosen above give
\[
 \widetilde L_{3,v_+}(x)=a-q\xi,\qquad
 \widetilde L_{3,v_-}(x)=a-q\xi'.
\]
Their corresponding local absolute values are both
$|a-q\xi|^{1/2}$, since $\tau(a-q\xi')=a-q\xi$.
Although $|a-q\xi'|$ can be large in the reference embedding,
this is not its absolute value at $v_-$.
The forms used at the two places are different, as permitted in the
definition of twisted height.
\end{remark}

All local determinants are $1$. The forms belong to the list
$X_1,X_2,X_3$ and the $d$ conjugates of
$-\xi X_1+\xi X_2+X_3$.
Consequently $\Delta_{\mathcal L}=1$ and $R=d+3$ is admissible;
the forms, field, and $H_{\mathcal L}$ are independent of $N,\eps,l$.
The identities \eqref{eq:rawsums} lift to
\[
 \sum_{v,i}\widetilde e_{iv}=-\delta,\qquad
 \sum_v\max_i\widetilde e_{iv}=1.
\]
Define
\begin{equation}\label{eq:center}
 \overline e_v=\frac13\sum_i\widetilde e_{iv},\qquad
 c_{iv}=\frac{\widetilde e_{iv}-\overline e_v}{1+\delta/3},
 \qquad Q_j=\Psi_j^{1+\delta/3}.
\end{equation}
The weights have finite support, sum to zero at each place, and obey
\[
 \sum_v\max_i c_{iv}
 =\frac{\sum_v\max_i\widetilde e_{iv}-\sum_v\overline e_v}
        {1+\delta/3}=1.
\]
Thus \eqref{eq:weightconditions} holds. By \eqref{eq:lift},
\[
 \|\widetilde L_{iv}(x_j)\|_v Q_j^{-c_{iv}}
 \le\Psi_j^{\overline e_v}.
\]
Taking maxima and multiplying over the places gives
\begin{equation}\label{eq:conversion}
 \HH_{\mathcal L,c,Q_j}(x_j)
 \le\Psi_j^{-\delta/3}
 =Q_j^{-\delta/(3+\delta)}
 =Q_j^{-\theta}.
\end{equation}
This is precisely the inequality required in
Corollary~\ref{thm:intervals}. Conjugating the forms in the passage to
$K$ is essential: the argument uses a small error at the specified
real embedding, without asserting the same error for every conjugate.

\section{Counting and proof of the main theorem}\label{sec:counting}

\begin{lemma}\label{lem:intervalcount}
Let $(Q_j)$ be a sequence of real numbers greater than $1$ such that
$Q_{j+1}>Q_j^2$. For $A>1$ and $\omega>1$, at most
$1+\lceil\log_2\omega\rceil$ of the parameters $Q_j$ belong to
$[A,A^\omega)$.
\end{lemma}
\begin{proof}
If $q$ such parameters occur, their logarithms increase by a factor
greater than $2^{q-1}$ from the first to the last. The interval
condition gives $2^{q-1}<\omega$ when $q\ge2$, which implies the
stated bound. The case $q\le1$ is immediate.
\end{proof}

\begin{proposition}\label{prop:count}
Assume \eqref{eq:LC}, with $F$ as in \eqref{eq:envelope},
$u>0$, and $\gamma\ge0$. Then, for every sufficiently large even $N$,
\begin{equation}\label{eq:count}
 \frac N2
 \ll_{b,\xi}\eps^{-3}(\log(2/\eps))^2
 \ll_{b,\xi,C,u,\gamma,k_0,\ell_1}
 N^{3u}(\log(2N))^{3u+2-3\gamma},
\end{equation}
where $\eps$ is defined by \eqref{eq:eps}.
\end{proposition}
\begin{proof}
Use the sequence of Lemma~\ref{lem:selection} and the parameters
\eqref{eq:eps}. By Proposition~\ref{prop:occupancy}, the exceptional
subspace supplied by Corollary~\ref{thm:intervals}, for any fixed
interval $I_l$, contains at most $B_{b,\xi}$ selected vectors.

The parameter sequence \eqref{eq:center} satisfies $Q_{j+1}>Q_j^2$.
By Lemma~\ref{lem:intervalcount}, each of the intervals supplied by
Corollary~\ref{thm:intervals} contains
$O_{b,\xi}(\log(2/\theta))$ selected parameters.
The parameters below $C_0$ can be discarded by restricting to the
indices $N/2\le j\le N$. Indeed,
\begin{equation}\label{eq:threshold}
 \theta^{-1}\ll\eps^{-1}\ll F(\ell_N)\ll(\log\ell_N)^u
             \ll(N\log(2N))^u,
 \qquad
 \log C_0=O_{b,\xi}(1)+O(\theta^{-1}),
\end{equation}
whereas $t_j\ge2^{j-1}t_1$ and
$\log Q_j=(1+\delta/3)t_j\log b$.
Thus $Q_j\ge C_0$ for all such indices when $N$ is large,
uniformly over the intervals $I_l$.

It follows from \eqref{eq:conversion} and
Corollary~\ref{thm:intervals} that each $I_l$ contains at most
\[
 B_{b,\xi}
 +O_{b,\xi}\bigl(\theta^{-2}(\log(2/\theta))^2\bigr)
\]
of these selected vectors. Summing over $k=O(\eps^{-1})$ intervals
and using \eqref{eq:exponentbounds}, we obtain
\begin{equation}\label{eq:countF}
 \frac N2
 \ll_{b,\xi}\eps^{-3}(\log(2/\eps))^2
 \ll_{b,\xi,C} F(\ell_N)^3(\log(2F(\ell_N)))^2.
\end{equation}

Put $y_N=\log\ell_N$. Then
$F(\ell_N)=y_N^u(\log y_N)^{-\gamma}$ and
$\log(2F(\ell_N))\ll_{u,\gamma}\log y_N$.
Consequently the right side of \eqref{eq:countF} is at most a constant
times
\[
 y_N^{3u}(\log y_N)^{2-3\gamma}.
\]
By \eqref{eq:loggrowth}, $y_N\ll N\log(2N)$ and
$\log y_N\asymp\log(2N)$. Since $3u>0$, this proves the second
inequality in \eqref{eq:count}. The two-sided comparison for
$\log y_N$ justifies the conclusion also when $2-3\gamma<0$.
\end{proof}

\begin{proof}[Proof of Theorem~\ref{thm:main}]
If the limsup in \eqref{eq:main} were finite for some
$\eta<1/3$, choose $\max(\eta,0)<u<1/3$.
The bound \eqref{eq:LC} would then hold with $\gamma=0$ and suitable
fixed $C,k_0$. Proposition~\ref{prop:count} gives
$N\ll N^{3u}(\log(2N))^{3u+2}=o(N)$, a contradiction.
\end{proof}

\begin{proof}[Proof of Theorem~\ref{thm:endpoint}]
If the limsup in \eqref{eq:endpoint} were finite for some fixed
$\gamma>1$, then \eqref{eq:LC} would hold with $u=1/3$ and that value
of $\gamma$. Proposition~\ref{prop:count} would give
\[
 N\ll N(\log(2N))^{3-3\gamma}=o(N),
\]
again a contradiction.
\end{proof}

\section{The role of fixed weights}\label{sec:scope}

The factor $\eps^{-1}$ from the partition of $r/t$ is relevant to
the exponent in \eqref{eq:count}. The following elementary observation
specifies the restriction imposed by the coordinatewise estimates
used here.

\begin{proposition}\label{prop:bins}
Fix $\eps>0$. Consider exponent vectors obtained by uniform
coordinatewise majorisation of
\[
 (1,\rho,-\eps)\quad\text{at infinity},\qquad
 (-w_p,-w_p\rho,0)\quad\text{at }p\mid b,
\]
for all $\rho\in[A,B]\subset[0,1]$.
Every such fixed exponent system has total sum at least
$-\eps+B-A$. To retain a negative total sum of magnitude at least
$\eps/2$, it is necessary that $B-A\le\eps/2$.
\end{proposition}
\begin{proof}
The second coordinate requires an exponent at least $B$ at infinity
and at least $-w_pA$ at $p\mid b$. The other coordinates require
the exponents already displayed. Their total is therefore at least
\[
 1+B-\eps-\sum_{p\mid b}w_p(1+A)=-\eps+B-A.
 \qedhere
\]
\end{proof}

This proposition concerns the displayed coordinatewise bounds;
it does not exclude estimates that exploit additional relations
among the points, or a theorem permitting weights to vary.
The parameter theorem used here keeps the weights fixed.

At $u=1/3$ and $\gamma=0$, \eqref{eq:count} becomes
$N\ll N(\log(2N))^3$, which is compatible with arbitrarily large $N$.
Theorem~\ref{thm:endpoint} instead uses $\gamma>1$.
At $\gamma=1$ the count gives only $N\ll N$; neither this borderline
case nor the unweighted endpoint is obtained.
The proof also extracts no effective upper bound for a length
witnessing \eqref{eq:main} or \eqref{eq:endpoint}. Its conclusion leaves open both
polynomial improvements in $L$ and positive entropy.

\appendix
\section{Comparison with covering theorems}\label{app:covering}

The same sequence and uniform subspace bound can be combined with
two covering theorems. We record the calculation to distinguish
the contribution of the subspace bound from that of the interval
theorem. These are weaker consequences of the same argument,
not additional assumptions in Theorem~\ref{thm:main}.

First, \cite[Theorem~5.1]{BE08} applies directly to the unnormalised
system \eqref{eq:rawineq}. Throughout this comparison take $\gamma=0$,
so that $F(x)=(\log x)^u$. In dimension $3$, for fixed forms of bounded
degree and height, its number of exceptional rational subspaces is
$O_{b,\xi}((1+\delta^{-1})^7)$.
Its required size threshold is
\[
 \Psi>\max(2H_*,3^{6/\delta}),
\]
where $H_*$ is a fixed bound for the heights $H(1,\alpha_1,\alpha_2,\alpha_3)$
of their coefficient vectors $(\alpha_1,\alpha_2,\alpha_3)$.
The hypotheses on exponents are $e_{i,\infty}\le1$,
$e_{i,p}\le0$ at finite places, and total sum $-\delta<0$.
They all hold in \eqref{eq:raw}. Here there are four distinct forms,
each of degree at most $d$, and each local determinant is $1$.
The last half of the selected vectors exceeds the threshold by the
same comparison as in \eqref{eq:threshold}.
Proposition~\ref{prop:occupancy} therefore gives
$O_{b,\xi}(\eps^{-7})$ points per interval $I_l$ and
$O_{b,\xi}(\eps^{-8})$ in total. Under \eqref{eq:LC}, this is
\[
 N\ll(N\log(2N))^{8u},
\]
a contradiction for $u<1/8$.
This combines the covering theorem of \cite{BE08} with
Proposition~\ref{prop:occupancy} of the present paper; the complexity
result stated in \cite[Theorem~2.1]{BE08} itself has range $\eta<1/11$.

Second, \cite[Theorem~2.1]{EF13} applies to the twisted-height system
of Section~\ref{sec:conversion}. For fixed $R$, it gives
$O_R(\theta^{-3}(\log(2/\theta))^2)$ proper subspaces covering all
solutions with $Q\ge C_0$; more precisely, for each such parameter,
the corresponding solutions are contained in one member of this
fixed finite collection. All hypotheses and the threshold have
already been checked.
The resulting total is
$O_{b,\xi}(\eps^{-4}(\log(2/\eps))^2)$, and hence
\[
 N\ll N^{4u}(\log(2N))^{4u+2},
\]
which is impossible for $u<1/4$.

For comparison, the interval theorem used in the main proof gives
the following counts. Constants in the table depend on the fixed
$b,\xi$.
\begin{center}
\small
\begin{tabular}{@{}lll@{}}
\toprule
Input & Total selected-point bound & Range\\
\midrule
\cite[Theorem~5.1]{BE08}
 & $O(\eps^{-8})$ & $\eta<1/8$\\
\cite[Theorem~2.1]{EF13}
 & $O(\eps^{-4}(\log(2/\eps))^2)$ & $\eta<1/4$\\
\cite[Theorem~2.3]{EF13}
 & $O(\eps^{-3}(\log(2/\eps))^2)$ & $\eta<1/3$\\
\bottomrule
\end{tabular}
\end{center}
The theorem numbers for \cite{EF13} refer to the published version.

\begin{thebibliography}{9}
\bibitem{AB07}
B.~Adamczewski and Y.~Bugeaud,
On the complexity of algebraic numbers I. Expansions in integer bases,
\emph{Ann. of Math.} (2) \textbf{165} (2007), no.~2, 547--565.
\href{https://doi.org/10.4007/annals.2007.165.547}
     {doi:10.4007/annals.2007.165.547}.

\bibitem{BE08}
Y.~Bugeaud and J.-H.~Evertse,
On two notions of complexity of algebraic numbers,
\emph{Acta Arith.} \textbf{133} (2008), no.~3, 221--250.
\href{https://doi.org/10.4064/aa133-3-3}
     {doi:10.4064/aa133-3-3}.

\bibitem{EF13}
J.-H.~Evertse and R.~G.~Ferretti,
A further improvement of the quantitative Subspace Theorem,
\emph{Ann. of Math.} (2) \textbf{177} (2013), no.~2, 513--590.
\href{https://doi.org/10.4007/annals.2013.177.2.4}
     {doi:10.4007/annals.2013.177.2.4}.
\end{thebibliography}
\end{document}
